% ============================================================================
%  Chapter 6 -- Number Representation: Real Numbers              (Lecture 06)
% ============================================================================
\chapter{Number Representation: Real Numbers}
\label{ch:reals}

\epigraphline{I may count for little, it's true, said the One to the Zero, but
what are you worth? Nothing: nothing at all.}{Trilussa, \emph{Nummeri}, 1944}

\begin{chapgoals}
By the end of this chapter you will be able to
\begin{itemize}
\item encode signed integers in one's complement and in excess-$k$, and say what each is good for;
\item represent a fractional quantity in fixed point and compute with it using integer instructions;
\item explain mantissa, exponent and normalisation in floating point;
\item encode and decode a number in IEEE~754 single precision;
\item recognise zero, denormalised numbers, infinity and NaN;
\item say why the ESP32-C6 has no floating-point instructions and what \prg games do instead.
\end{itemize}
\end{chapgoals}

% ----------------------------------------------------------------------------
\section{Warm-up on the machine}

A ball must cross the screen slowly. One pixel per frame is too fast; zero
pixels per frame is not a movement. The ball needs something in between:
0.375 pixels per frame. But the processor of the kit is RV32IMAC: it has
integer instructions only, and no floating point at all.
Listing~\ref{lst:real:ball} solves the problem with a trick.

\asmfile[caption={Three-eighths of a pixel per frame, with integers only.},label=lst:real:ball]{code/l06/s04.S}

The position is stored \emph{multiplied by 256}, and is divided by 256 only
when the ball is drawn. In units of $1/256$ of a pixel the speed is
$0.375\times 256 = 96$, an honest integer.

\begin{prgtry}
Run the program with \code{SPEED\_Q8} equal to 96, then 16, then 1.
\end{prgtry}

The trick has a name, \term{fixed point}. This chapter says where it comes
from and when it is not enough, in which case the general answer,
\term{floating point}, takes over. Before either, two more encodings of the
signed integers are needed, because floating point is built from them.

% ----------------------------------------------------------------------------
\section{More encodings for signed integers}

\subsection{One's complement}

\begin{definition}{One's complement}
On $n$ bits, a non-negative number $0\le x<2^{n-1}$ is represented by itself,
$X=x$; a negative number is represented by $X = 2^{n}-1-|x|$.
\end{definition}

The range goes from $-(2^{n-1}-1)$ to $+(2^{n-1}-1)$. Zero has two
representations, \texttt{0000} and \texttt{1111}, so that $2^{n}-1$ numbers
are represented, as in sign and magnitude. The reason for using it is that
negation is trivial: since $2^{n}-1$ is a string of ones, subtracting from it
complements every bit,
\[ -x \;=\; \mathtt{not}\ x. \]
Thus $6=\texttt{0110}$ and $-6=\texttt{1001}$; $5=\texttt{0101}$ and
$-5=\texttt{1010}$; $1=\texttt{0001}$ and $-1=\texttt{1110}$. The bitwise
\ins{not} of \riscv computes exactly this.

Arithmetic in one's complement works modulo $2^{n}-1$ rather than modulo
$2^{n}$, and this costs a correction.

\begin{center}\ttfamily
\begin{tabular}{@{}r@{\quad}r@{}}
\multicolumn{2}{c}{\textsf{\small No correction}}\\
3 + & 0011 +\\ -4 = & 1011 =\\\hline -1 & 1110\\ &
\end{tabular}\qquad\qquad
\begin{tabular}{@{}r@{\quad}r@{}}
\multicolumn{2}{c}{\textsf{\small Two negatives}}\\
-2 + & 1101 +\\ -3 = & 1100 =\\\hline & \textcolor{prgred}{1}\,1001 + 1 =\\ -5 & 1010
\end{tabular}\qquad\qquad
\begin{tabular}{@{}r@{\quad}r@{}}
\multicolumn{2}{c}{\textsf{\small Positive result}}\\
5 + & 0101 +\\ -2 = & 1101 =\\\hline & \textcolor{prgred}{1}\,0010 + 1 =\\ 3 & 0011
\end{tabular}
\end{center}

Whenever there is a carry out of the most significant bit, 1 must be added to
the result: the \term{end-around carry}. The correction is needed when both
operands are negative, or when the operands have opposite signs and the sum is
positive. An adder and a complementer are still enough; but the extra step
and the double zero made two's complement win.

\subsection{Excess-$k$ (biased) representation}

\begin{definition}{Excess-$k$}
In the \term{excess-$k$} representation every number is shifted upward by a
constant $k$, the \term{bias}: $X = x+k$. Typically $k=2^{n-1}$.
\end{definition}

On four bits with $k=8$ the pattern \texttt{0000} means $-8$, \texttt{1000}
means 0 and \texttt{1111} means $+7$. The decisive property is that the order
is preserved: if $x_1>x_2$ then $X_1>X_2$, so that signed numbers are compared
as if they were naturals. The price is paid in arithmetic: a sum must be
corrected by subtracting $k$, since $X_1+X_2 = (x_1+x_2)+2k$, and
multiplication and division are much more awkward than with the complements.
Excess-$k$ is therefore preferred where only additions and comparisons are
needed, and that is exactly the situation of the exponent of a floating-point
number.

\subsection{Four bits, five readings}

Table~\ref{tab:real:five} completes the table of Chapter~\ref{ch:integers}.
Sign and magnitude and one's complement have two zeros. Two's complement and
excess-8 cover the same range, $-8$ to $+7$, with a single zero; they differ
only in the most significant bit, which is inverted. Only excess-$k$ keeps
the patterns in the same order as the numbers.

\begin{table}[htb]
\centering\small
\caption{Four bits, five readings.}\label{tab:real:five}
\setlength{\tabcolsep}{4.5pt}
\begin{tabular}{@{}cccccc@{\qquad}cccccc@{}}
\toprule
\textbf{Bits} & \textbf{Uns.} & \textbf{S-M} & \textbf{1's} & \textbf{2's} & \textbf{Exc-8} &
\textbf{Bits} & \textbf{Uns.} & \textbf{S-M} & \textbf{1's} & \textbf{2's} & \textbf{Exc-8}\\
\midrule
\texttt{0000} & 0 & $+0$ & $+0$ & 0 & $-8$ & \texttt{1000} & 8 & $-0$ & $-7$ & $-8$ & 0\\
\texttt{0001} & 1 & $+1$ & $+1$ & $+1$ & $-7$ & \texttt{1001} & 9 & $-1$ & $-6$ & $-7$ & $+1$\\
\texttt{0010} & 2 & $+2$ & $+2$ & $+2$ & $-6$ & \texttt{1010} & 10 & $-2$ & $-5$ & $-6$ & $+2$\\
\texttt{0011} & 3 & $+3$ & $+3$ & $+3$ & $-5$ & \texttt{1011} & 11 & $-3$ & $-4$ & $-5$ & $+3$\\
\texttt{0100} & 4 & $+4$ & $+4$ & $+4$ & $-4$ & \texttt{1100} & 12 & $-4$ & $-3$ & $-4$ & $+4$\\
\texttt{0101} & 5 & $+5$ & $+5$ & $+5$ & $-3$ & \texttt{1101} & 13 & $-5$ & $-2$ & $-3$ & $+5$\\
\texttt{0110} & 6 & $+6$ & $+6$ & $+6$ & $-2$ & \texttt{1110} & 14 & $-6$ & $-1$ & $-2$ & $+6$\\
\texttt{0111} & 7 & $+7$ & $+7$ & $+7$ & $-1$ & \texttt{1111} & 15 & $-7$ & $-0$ & $-1$ & $+7$\\
\bottomrule
\end{tabular}
\end{table}

% ----------------------------------------------------------------------------
\section{Fixed point}

\subsection{Numbers in fixed point}

In fixed point the integer part and the fractional part of a number are stored
side by side in one word. The position of the point is fixed and implicit: it
is not stored, the program knows it.

\begin{definition}{Fixed point}
With $f$ fractional bits a number $x$ is stored as the integer
$X = x\cdot 2^{f}$ (rounded). The \term{resolution} is $2^{-f}$. The notation
Q$m.f$ denotes a format with $m$ integer bits and $f$ fractional bits.
\end{definition}

Usually the whole word is in two's complement. Fixed point is how games,
audio and control software handle fractions on processors without a
floating-point unit. It is also how money should be stored: as an integer
number of cents.

\subsection{The ball, bit by bit: Q24.8}

The ball of the warm-up is in the format Q24.8
(Figure~\ref{fig:real:q248}). When its variable holds
$\hex{00002A60}=10848$, the ball is at $10848/256 = 42.375$ pixels: the upper
24 bits, \hex{00002A}, are the 42 whole pixels and the lower 8, \hex{60}, are
$96/256 = 0.375$.

\begin{figure}[htb]
\centering
\begin{bytefield}[bitwidth=0.029\textwidth,bitheight=8mm]{32}
\bitheader[endianness=big]{0,7,8,31}\\
\bitbox{24}[bgcolor=prgaccent!20]{\texttt{0000 0000 0000 0000 0010 1010}} & \bitbox{8}[bgcolor=prgamber!30]{\texttt{0110 0000}}\\
\bitbox[]{24}{\footnotesize integer part: 24 bits\\\texttt{0x00002A} = 42 pixels} & \bitbox[]{8}{\footnotesize fraction: 8 bits\\\texttt{0x60} = 96/256}
\end{bytefield}
\caption{\code{ball\_x\_q8 = 0x00002A60}: the ball is at 42.375 pixels.}
\label{fig:real:q248}
\end{figure}

After one more frame the value is \hex{2AC0} (42.75), then \hex{2B20}
(43.125): the ball crosses the boundary of a pixel every two or three frames,
and on average advances by 0.375.

Arithmetic in fixed point is integer arithmetic with a little care
(Table~\ref{tab:real:ops}). Sums and differences need none, because the
points of the two operands are aligned. A product of two numbers scaled by
$2^{f}$ is scaled by $2^{2f}$, and must be shifted back by $f$. The shift
that returns to whole pixels is the \emph{arithmetic} one, \ins{srai}, which
keeps the sign (Chapter~\ref{ch:integers}).

\begin{table}[htb]
\centering\small
\caption{Operations in fixed point.}\label{tab:real:ops}
\begin{tabular}{@{}lll@{}}
\toprule
\textbf{Operation} & \textbf{In fixed point} & \textbf{\riscv}\\
\midrule
add, subtract & as integers: the points are aligned & \ins{add}, \ins{sub}, \ins{addi}\\
to whole pixels & shift right by $f$, keeping the sign & \code{srai a0, t1, 8}\\
multiply & multiply, then shift right by $f$ & \ins{mul}, then \ins{srai}\\
\bottomrule
\end{tabular}
\end{table}

\subsection{Rational numbers: the approximation error}

The real numbers are dense: even in a limited interval, not every number can
be represented with a finite number of bits. Each number is therefore
associated with the nearest representable one. In Q24.8 the number 0.1
becomes $26/256 = 0.1015625$, because $0.1\times 256 = 25.6$ is rounded to 26:
an error of about 1.6\%.

The approximation error is at most half the resolution: $1/512$ with eight
fractional bits. \term{Underflow} is the case of a number close to zero, but
not zero, that is represented by zero: with a speed of 0.001 pixels per frame
the constant would be rounded to 0, and the ball would never move.

\begin{prgconcept}
In fixed point the \emph{absolute} error is the same everywhere on the number
line. The resolution of eight fractional bits is $1/256$: nothing smaller can
be told apart from zero.
\end{prgconcept}

% ----------------------------------------------------------------------------
\section{Floating point}

Fixed point spends all its bits on one fixed range. Floating point spends
some of them to say \emph{where} the range is: it trades range for precision.

\subsection{Real numbers in floating point}

\begin{definition}{Floating point}
A real number $x$ is represented in floating point by a pair $(m,e)$ such that
\[ x = m\cdot b^{\,e}, \]
where $m$ is the \term{mantissa}, $e$ the \term{exponent} and $b$ the base.
\end{definition}

The point ``floats'' because the exponent says where it is. This is the
scientific notation of the physics laboratory. Infinitely many pairs
represent the same number,
\[
346.09801 \;=\; (346.09801,\,0) \;=\; (346098.01,\,-3) \;=\; (0.034609801,\,4),
\]
and \term{normalisation} chooses a single representative: the point goes
immediately after the first significant digit, $(3.4609801,\,2)$.

\subsection{The range of a small decimal format}

Consider a format in base 10 with four digits for the mantissa and two for
the exponent (Figure~\ref{fig:real:line}). It represents the numbers from
$-9.999\times 10^{99}$ to $-1.000\times 10^{-99}$, zero, and the numbers from
$+1.000\times 10^{-99}$ to $+9.999\times 10^{99}$. An \term{overflow} occurs
when the magnitude is too large for the exponent; an \term{underflow} when it
is too small, and the number collapses to zero.

\begin{figure}[htb]
\centering
\begin{tikzpicture}[x=1cm,y=1cm]
  \draw[<->,thick,prgink] (-7,0) -- (7,0);
  \fill[prgaccent!30] (-5.6,-0.22) rectangle (-1.0,0.22);
  \fill[prgaccent!30] (1.0,-0.22) rectangle (5.6,0.22);
  \fill[prgred!25] (-7,-0.22) rectangle (-5.6,0.22);
  \fill[prgred!25] (5.6,-0.22) rectangle (7,0.22);
  \fill[prgamber!35] (-1.0,-0.22) rectangle (1.0,0.22);
  \draw[thick] (0,-0.3) -- (0,0.3) node[above,lbl]{0};
  \foreach \x/\l in {-5.6/{$-9.999\times10^{99}$},-1.0/{$-1.000\times10^{-99}$},1.0/{$1.000\times10^{-99}$},5.6/{$9.999\times10^{99}$}}
    \draw (\x,-0.3) -- (\x,0.3) node[above,font=\scriptsize]{\l};
  \node[lbl,text=prgred] at (-6.3,-0.6) {overflow};
  \node[lbl,text=prgred] at (6.3,-0.6) {overflow};
  \node[lbl,text=prgamber] at (0,-0.6) {underflow};
  \node[lbl,text=prgaccent] at (-3.3,-0.6) {representable};
  \node[lbl,text=prgaccent] at (3.3,-0.6) {representable};
\end{tikzpicture}
\caption{A decimal format with four digits of mantissa and two of exponent
(not to scale).}
\label{fig:real:line}
\end{figure}

\subsection{Approximation: absolute and relative error}

In such a format the absolute error is not constant. Near zero the distance
between two consecutive numbers is
$1.001\times 10^{-99}-1.000\times 10^{-99} = 0.001\times 10^{-99}$; near the
upper limit it is
$9.999\times 10^{99}-9.998\times 10^{99} = 0.001\times 10^{99}$. The absolute
error is small on small numbers and large on large numbers. What is constant
along the whole axis is the \emph{relative} error, because the number of
significant digits, four, is the same everywhere.

\begin{prgconcept}
Fixed point: constant absolute error, relative error that explodes near zero.
Floating point: constant relative error, absolute error that grows with the
number.
\end{prgconcept}

\subsection{IEEE 754 single precision}

Until 1985 every manufacturer had its own floating-point format, and programs
gave different results on different machines. The standard IEEE~754 (1985,
revised in 2008 and 2019) put an end to this \citep{ieee754}. Its
single-precision format divides 32 bits into three fields
(Figure~\ref{fig:real:ieee}).

\begin{figure}[htb]
\centering
\begin{bytefield}[bitwidth=0.029\textwidth,bitheight=8mm]{32}
\bitheader[endianness=big]{0,22,23,30,31}\\
\bitbox{1}[bgcolor=prgred!30]{S} & \bitbox{8}[bgcolor=prgamber!30]{E} & \bitbox{23}[bgcolor=prgaccent!20]{F}\\
\bitbox[]{1}{\footnotesize sign} & \bitbox[]{8}{\footnotesize exponent: 8 bits\\excess-127} & \bitbox[]{23}{\footnotesize fraction: 23 bits\\magnitude, after the hidden 1}
\end{bytefield}
\[ \text{value} \;=\; (-1)^{S}\times 1.F\times 2^{\,E-127} \]
\caption{IEEE 754 single precision.}
\label{fig:real:ieee}
\end{figure}

\begin{itemize}
\item The mantissa is in sign and magnitude: one sign bit and 23 bits of
  magnitude.
\item A normalised binary mantissa always starts with 1. That bit is
  \term{hidden}: it is not stored, and one bit of precision is gained.
\item The exponent is in excess-127 (not 128): the true exponents go from
  $-126$ to $+127$. Thanks to the excess encoding, and to the position of the
  exponent before the fraction, two positive floating-point numbers can be
  compared as if they were naturals.
\item The patterns $E=0$ and $E=255$ are reserved: zero and denormalised
  numbers, infinity and NaN.
\end{itemize}

\begin{example}{$-5.125$ in IEEE 754}
\begin{enumerate}[nosep]
\item Convert the two parts: $5=101_2$ and $0.125=1/8=0.001_2$, so
  $5.125=101.001_2$.
\item Normalise: $101.001_2 = 1.01001_2\times 2^{2}$.
\item Sign $S=1$; exponent $E = 2+127 = 129 = \texttt{1000\,0001}$.
\item Fraction $F=\texttt{01001}$ followed by zeros: the leading 1 is hidden.
\end{enumerate}
\[
\underbrace{\texttt{1}}_{S}\;\underbrace{\texttt{1000\,0001}}_{E}\;
\underbrace{\texttt{0100\,1000\,0000\,0000\,0000\,000}}_{F}
\]
Regrouped in nibbles, \texttt{1100 0000 1010 0100 0000 \ldots} is
\hex{C0A40000}.
\end{example}

\subsection{IEEE 754 at a glance}

Table~\ref{tab:real:formats} compares the two formats in common use, and
Table~\ref{tab:real:special} lists the meaning of all the patterns.

\begin{table}[htb]
\centering\small
\caption{IEEE 754 single and double precision.}\label{tab:real:formats}
\begin{tabular}{@{}lll@{}}
\toprule
 & \textbf{Single precision} & \textbf{Double precision}\\
\midrule
Total bits & 32 & 64\\
Sign bits & 1 & 1\\
Exponent bits & 8 & 11\\
Fraction bits & 23 & 52\\
Exponent encoding & excess-127 & excess-1023\\
Exponent range & $-126\ldots+127$ & $-1022\ldots+1023$\\
Smallest normalised & $2^{-126}\approx 1.2\times 10^{-38}$ & $2^{-1022}\approx 2.2\times 10^{-308}$\\
Largest & about $2^{128}\approx 3.4\times 10^{38}$ & about $2^{1024}\approx 1.8\times 10^{308}$\\
Smallest denormalised & about $1.4\times 10^{-45}$ & about $4.9\times 10^{-324}$\\
Decimal digits & about 7 & about 16\\
\bottomrule
\end{tabular}
\end{table}

\begin{table}[htb]
\centering\small
\caption{IEEE 754: what the patterns mean.}\label{tab:real:special}
\begin{tabular}{@{}lllll@{}}
\toprule
\textbf{Kind} & \textbf{Sign} & \textbf{Exponent $E$} & \textbf{Fraction $F$} & \textbf{Value}\\
\midrule
Normalised & 0 or 1 & $0<E<\max$ & any pattern & $(-1)^{S}\times 1.F\times 2^{E-\mathit{bias}}$\\
Denormalised & 0 or 1 & all zeros & not zero & $(-1)^{S}\times 0.F\times 2^{1-\mathit{bias}}$\\
Zero & 0 or 1 & all zeros & all zeros & $+0$ or $-0$\\
Infinity & 0 or 1 & all ones & all zeros & $+\infty$ or $-\infty$\\
Not a Number & 0 or 1 & all ones & not zero & NaN: an invalid result\\
\bottomrule
\end{tabular}
\end{table}

\term{Denormalised} (or subnormal) numbers fill the gap of underflow around
zero with a precision that decreases gradually. The division $1/0$ gives
infinity; $0/0$ and the square root of $-1$ give NaN. NaN is not equal to
anything, itself included.

\subsection{Floating point is not the real numbers}

A floating-point number is a rational number chosen from a finite set, and
arithmetic on such numbers obeys laws that are not those of the real field
\citep{goldberg1991}.

\begin{itemize}
\item 0.1 has no finite binary expansion: it is stored approximately, like 0.8
  in Chapter~\ref{ch:integers}.
\item Two floating-point results should never be tested for equality: their
  difference is compared with a tolerance (Listing~\ref{lst:real:c}).
\item Addition is not associative: $(a+b)+c$ can differ from $a+(b+c)$.
\item Adding a very small number to a very large one can leave the large one
  unchanged.
\item For money and for the positions of a game, scaled integers are often the
  right tool.
\end{itemize}

\cfile[caption={The most famous sum in computing.},label=lst:real:c,numbers=none]{code/l06/s22.c}

\subsection{Beyond single and double}

The choice of a representation is an engineering decision that is still being
made today. \emph{Half precision}, added to the standard in 2008, has sixteen
bits: 1 of sign, 5 of exponent, 10 of fraction; it gives about three decimal
digits in half the memory of single precision. It matters because neural
networks hold thousands of millions of numbers: halving each one halves
memory, bandwidth and energy. Many models go further and are stored, by
\term{quantisation}, as 8-bit integers with a scale factor: fixed point again,
at industrial scale. Bits are a budget; range, precision and cost pull in
different directions, on a GPU as on a Q24.8 ball.

% ----------------------------------------------------------------------------
\section{Back to the machine}

\subsection{Floating point on RISC-V, and on the kit}

\riscv is modular: a small mandatory base plus optional extensions
(Table~\ref{tab:real:rv}). A \riscv processor with the F and D extensions has
thirty-two more registers, \reg{f0} to \reg{f31}, for floating-point values,
and instructions such as \ins{fadd.s}, \ins{fmul.s}, \ins{flw} and \ins{fsw}
that work on IEEE~754 numbers; \ins{fcvt.w.s} converts to an integer and
\ins{fcvt.s.w} from an integer. Such processors are found in laptops, servers
and supercomputers.

The ESP32-C6 is RV32IMAC: it has no F and no D, and therefore no
floating-point instructions at all. A C program can still use the type
\code{float}: the compiler calls software routines such as
\code{\_\_addsf3}, which implement exactly the algorithm of this chapter with
shifts, masks and integer additions, at a cost of tens to hundreds of integer
instructions per operation. \prg games use integers and fixed point: fast,
exact enough, and visible in the debugger.

\begin{table}[htb]
\centering\small
\caption{Reading the name RV32IMAC.}\label{tab:real:rv}
\begin{tabularx}{\textwidth}{@{}llLc@{}}
\toprule
\textbf{Letter} & \textbf{Extension} & \textbf{What it adds} & \textbf{ESP32-C6}\\
\midrule
RV32I & base integer set & 32-bit registers, add, logic, load and store, branches & yes\\
M & multiply and divide & \ins{mul}, \ins{mulh}, \ins{div}, \ins{rem} and their unsigned forms & yes\\
A & atomic & instructions for safe sharing between tasks & yes\\
C & compressed & 16-bit encodings of the most common instructions & yes\\
F & single precision & \reg{f0} to \reg{f31} and IEEE 754 single-precision instructions & no\\
D & double precision & IEEE 754 double-precision instructions & no\\
\bottomrule
\end{tabularx}
\end{table}

A microcontroller picks the extensions that it needs. This is one reason why
the same architecture scales from a chip that costs two euros to a
supercomputer.

% ----------------------------------------------------------------------------
\section{Check your understanding}

\begin{quiz}{Which signed encoding keeps the bit patterns in the same order as the numbers?}
\item Sign and magnitude
\item One's complement
\item Two's complement
\item Excess-$k$
\end{quiz}

\begin{quiz}{In Q24.8 fixed point, which real number does the word \hex{00000180} represent?}
\item 384
\item 1.5
\item 1.8
\item 0.375
\end{quiz}

\begin{quiz}{Which real number is the IEEE 754 single-precision pattern \hex{40400000}?}
\item 2.0
\item 3.0
\item 4.0
\item 0.75
\end{quiz}

\begin{quiz}{In IEEE 754, an exponent field of all ones with a fraction of all zeros means:}
\item zero
\item a denormalised number
\item infinity
\item NaN
\end{quiz}

\section{Exercises}

\begin{exercise}[paper first, then \prg]
\begin{enumerate}[nosep]
\item Write $-6$ on four bits in sign and magnitude, one's complement, two's
  complement and excess-8.
\item Encode 10.5 and 0.15625 in IEEE 754 single precision; give the result in
  hexadecimal.
\item Decode \hex{BF800000}.
\item Run the ball with \code{SPEED\_Q8} equal to 96, 16 and 1. After how many
  frames does it move by one pixel in each case?
\item Give the ball a vertical speed too, and make it bounce: negate the speed
  with \ins{neg} when it reaches an edge.
\end{enumerate}
\expected{On paper: \texttt{1110}, \texttt{1001}, \texttt{1010},
\texttt{0010}; \hex{41280000} and \hex{3E200000}; $-1.0$. On \prg: the ball
advances one pixel every 2.67, 16 and 256 frames; with the bounce it moves
diagonally and stays on the screen.}
\end{exercise}

\section{Glossary}

\begin{glossarytable}
one's complement & complemento a uno (alla base diminuita)\\
excess-$k$ (biased) & eccesso-$k$\\
fixed point & virgola fissa\\
floating point & virgola mobile\\
single / double precision & singola / doppia precisione\\
rounding & arrotondamento\\
mantissa / fraction & mantissa\\
exponent & esponente\\
normalisation & normalizzazione\\
hidden bit & bit nascosto\\
denormalised & denormalizzato\\
Not a Number (NaN) & non un numero (NaN)\\
\end{glossarytable}

\section{Conclusions and reading}

Integers have several representations: 8, 16, 32 or 64 bits, signed or
unsigned, complements or excess. Fixed point stores a fraction as a scaled
integer: constant absolute error, the speed of integers. Floating point
stores a mantissa and an exponent: a huge range, constant relative error.
IEEE~754 is the universal standard for floating point, with special patterns
for zero, infinity and NaN. On a processor without the F extension, fixed
point is the natural choice.

\paragraph{Reading.} \citet{patterson2020}, section~3.5;
\citet{goldberg1991}; \emph{Playing with RISC-V}, the appendix on the runtime
interface, for playfields and parallax in Q8 \citep{prg32textbook}.

The first block is complete. The next chapter puts its four lectures to work
in a real program: the first part of \vermone.
